\documentclass[preview]{standalone}

\usepackage{amsmath}
\usepackage{amssymb}
\usepackage{stellar}
\usepackage{definitions}

\begin{document}

\id{la-scienza-cognitiva}
\genpage

\section{La scienza cognitiva}

\begin{snippetdefinition}{scienza-cognitiva-definition}{Scienza cognitiva}
    La \emph{scienza cognitiva} è un ambito interdisciplinare nato negli anni
    Settanta per studiare la cognizione integrando psicologia, informatica,
    filosofia, linguistica, neuroscienze e modelli computazionali.
\end{snippetdefinition}

\begin{snippet}{paradigmi-scienza-cognitiva}
    I due paradigmi principali della scienza cognitiva sono:
    \begin{itemize}
        \item il \emph{modularismo}, secondo cui la mente è organizzata in
        moduli specializzati;
        \item il \emph{connessionismo}, che mette in relazione l'architettura
        biologica del cervello con l'architettura funzionale dell'attività
        cognitiva.
    \end{itemize}
    I modelli della mente possono essere distinti in modelli
    \emph{computer-style}, basati sull'analogia con il computer, e modelli
    \emph{brain-style}, basati sul funzionamento del cervello.
\end{snippet}

\section{Modularismo}

\begin{snippetdefinition}{modulo-cognitivo-definition}{Modulo}
    Un \emph{modulo} è una struttura specializzata della mente che funziona
    come sistema esperto in uno specifico ambito dell'interazione con
    l'ambiente.
\end{snippetdefinition}

\begin{snippet}{fodor-modularismo}
    Secondo Fodor, la mente è organizzata in moduli biologicamente determinati
    e specie-specifici. Il sistema cognitivo comprende trasduttori, sistemi di
    input e sistemi centrali. I trasduttori registrano informazioni
    ambientali; i sistemi di input elaborano informazioni percettive e
    linguistiche; i sistemi centrali integrano le informazioni prodotte dai
    diversi sistemi di input.
\end{snippet}

\begin{snippetdefinition}{incapsulamento-moduli-definition}{Incapsulamento}
    L'\emph{incapsulamento} è la proprietà dei sistemi di input per cui essi
    non sono influenzati dagli altri sistemi di input né dai processi centrali.
    I processi centrali, invece, non sono incapsulati e interagiscono tra loro.
\end{snippetdefinition}

\section{Connessionismo}

\begin{snippetdefinition}{connessionismo-definition}{Connessionismo}
    Il \emph{connessionismo} è una prospettiva che spiega la cognizione tramite
    reti di molte unità interconnesse, ispirate alla struttura neurale del
    cervello. L'elaborazione dell'informazione avviene in parallelo e la
    conoscenza dipende da meccanismi di associazione.
\end{snippetdefinition}

\begin{snippet}{reti-neurali-artificiali}
    Il connessionismo fa riferimento alle reti neurali artificiali, cioè
    simulazioni che riproducono in modo semplificato alcune proprietà del
    sistema nervoso. Questa impostazione sostiene una concezione dinamica della
    mente, capace di adattarsi alle condizioni del momento e di autocorreggersi.
\end{snippet}

\includesnpt[width=58\%,src=/snippet/static-psychology/connectionism-neural-networks.jpg]{centered-img}

\begin{snippetdefinition}{mente-situata-definition}{Mente situata}
    La \emph{mente situata} è una mente costantemente immersa in un contesto
    immediato, dal quale dipendono le sue attività cognitive.
\end{snippetdefinition}

\begin{snippetdefinition}{mente-simulativa-definition}{Mente simulativa}
    La \emph{mente simulativa} è una mente capace di elaborare rappresentazioni
    che permettono di riprodurre, prevedere e immaginare fenomeni, oggetti,
    eventi e mondi possibili.
\end{snippetdefinition}

\begin{snippet}{mente-radicata}
    Una mente situata è anche una mente radicata nel corpo, cioè una
    \emph{embodied mind}. Secondo questa prospettiva, la conoscenza è fondata
    nell'esperienza e deriva da informazioni sensoriali, immaginative,
    linguistiche, affettive e motorie.
\end{snippet}

\end{document}
